% ==============================================================================
% Section: A Price-Cap Lens on the Shapiro Settlement
% Status: REWRITTEN (2026-05-02) — replaces former \section{Model} as part of
%   laser-focus rewrite. SFE machinery moved to references; appendix removed.
% ==============================================================================

\section{A Price-Cap Lens on the Shapiro Settlement}
\label{sec:theory}
\label{subsec:theory_ceiling}%

Standard price-control theory delivers a sharp prediction for a binding
ceiling on a competitively cleared market. When the regulator imposes a
cap $\bar p_S$ below the unconstrained clearing price $p^*$, demand at
$\bar p_S$ exceeds the quantity sellers will provide; the gap between
$D(\bar p_S)$ and $S(\bar p_S)$ is the shortage.
\citet{Weitzman1974_prices} formalizes the choice between price and
quantity instruments under uncertainty, and
\citet{GlaeserLuttmer2003_rent} document a distinct allocative
consequence of binding ceilings: misallocation across heterogeneous
demanders, qualitatively different from the standard deadweight
triangle.

The Shapiro settlement applies this framework directly to PJM's Reliability
Pricing Model. The settlement modifies the Variable Resource Requirement
curve to be flat for all $p \geq \bar p_S = \$325$/MW-day UCAP and adds
a flat segment at $\underline p_S \approx \$175$/MW-day for low prices
(Section~\ref{subsec:inst_shapiro}). The capped auction therefore clears
whatever capacity sellers offer at or below $\bar p_S$. If aggregate
offers at the cap fall short of the reliability requirement $q_b$ (the
Point~(b) quantity in the VRR curve), the auction clears at
$(\bar p_S, Q_{\text{offered}})$ with $Q_{\text{offered}} < q_b$, and
two consequences are visible separately:
a truncated price signal on the upper segment, and---if supply is
tight---a quantity shortfall against the planning target. Which of these
channels was operative in 2026/27 and 2027/28 is the empirical question
of Sections~\ref{sec:cap_incidence} and~\ref{sec:case_studies}.

\label{subsec:theory_sfe}%
To distinguish a binding cap from a redundant one, the paper uses the
symmetric supply function equilibrium (SFE) of
\citet{Klemperer1989_sfe} as a characterization of the unconstrained
equilibrium price. The setup follows the calibrated-electricity-market
template of \citet{Green1992_british}: $K$ symmetric pivotal sellers
with constant marginal cost $c$ submit nondecreasing supply schedules
$S_i(\cdot)$ against the publicly known VRR demand curve $D(\cdot)$,
alongside a competitive fringe of size $\bar Q_f$ pricing at $c$.
Equilibrium is selected by the \citet{Holmberg2008_unique_sfe} boundary
condition, which requires each seller's supply function to attain its
capacity at the design-level Point~(a) cap, denoted $\bar p$ to
distinguish it from the lower settlement cap $\bar p_S$. Under symmetry the
equilibrium reduces to a scalar first-order ODE in the per-firm supply
function $s(p)$, integrated backward from $s(\bar p) = \bar q$;
the derivation is standard \citep{Klemperer1989_sfe, Green1992_british,
Rudkevich1998_poolco} and is not reproduced here.

The model's content for this paper is binary. The \emph{at-cap regime}
($p^* = \bar p$) obtains whenever rivals' capacity constraints
simultaneously bind, that is, when $K$ pivotal sellers cannot meet
residual demand at any interior price. The \emph{interior regime}
($p^* < \bar p$) holds otherwise. The Shapiro cap $\bar p_S < \bar p$
binds in the at-cap regime by construction; in the interior regime it
binds only if $\bar p_S < p^*$. Numerical solution of the symmetric ODE
for the calibrations of Section~\ref{sec:calibration} yields the at-cap
regime for $K \leq 3$ and the interior regime for $K \geq 4$; the
price-cost margin falls by roughly 20--25 percentage points across
this regime change.\footnote{For 2026/27 with $c = \$150$/MW-day, the
symmetric SFE clears at the cap ($p^* = \bar p = \$329$, $L = 0.54$)
under $K = 3$ and at $p^* \approx \$228$ ($L = 0.34$) under $K = 4$.
The qualitative pattern is the same for the other six calibration
years.}

PJM's auction concentration is consistent with the $K \leq 3$ regime.
The IMM reports $\mathrm{RSI}_3 < 1$ in every sample year
(Table~\ref{tab:tps_results}), meaning the three largest sellers are
collectively pivotal at the reliability requirement. A symmetric $K=3$
calibration is therefore the natural choice and is maintained throughout.
The SFE plays a focused role in what follows: it identifies which
auctions feature a binding settlement cap, and it supplies an
unconstrained-market counterfactual price against which the cap's
effect on cleared MW can be measured in
Section~\ref{sec:case_studies}. In the at-cap regime, the SFE matches
the observed clearing price by construction, not as an independent
empirical test: its content for this paper is the regime
classification itself, with the dollar value of $\bar p$ read directly
from PJM's planning parameters.
